\section*{Sets and Indices}

\begin{itemize}
    \item $M$: set of microcomputer models, with $M=\{A,B\}$.
    \item $P$: set of production processes, with $P=\{I,II\}$.
\end{itemize}

\section*{Parameters}

\begin{itemize}
    \item $h_{pm}$: processing time (hours per unit) required by model $m \in M$ on process $p \in P$.
    \item $C_p$: maximum weekly processing capacity (hours) available for process $p \in P$.
    \item $\pi_m$: profit per unit produced of model $m \in M$.
    \item $L_m$: minimum weekly production requirement (units) for model $m \in M$.
    \item $\Pi^{\min}$: minimum required weekly total profit.
\end{itemize}

Using the provided data, these values are:
\[
h_{IA}=4,\quad h_{IB}=6,\quad h_{II,A}=3,\quad h_{II,B}=2,
\]
\[
C_I=150,\quad C_{II}=70,
\]
\[
\pi_A=300,\quad \pi_B=450,
\]
\[
L_A=10,\quad L_B=15,\quad \Pi^{\min}=10000.
\]

\section*{Decision Variables}

\begin{itemize}
    \item $x_A \ge 0$ (code: \texttt{x\_A}): quantity of model A produced in the week.
    \item $x_B \ge 0$ (code: \texttt{x\_B}): quantity of model B produced in the week.
\end{itemize}

The implementation treats both variables as continuous, so fractional production quantities are allowed.

\section*{Objective}

Maximize total weekly profit:
\[
\max \; 300x_A + 450x_B.
\]

Equivalently, in indexed form:
\[
\max \; \sum_{m \in M} \pi_m x_m.
\]

\section*{Constraints}

Process capacity limits:
\[
4x_A + 6x_B \le 150 \qquad \text{(Process I)}
\]
\[
3x_A + 2x_B \le 70 \qquad \text{(Process II)}
\]

Minimum production requirements:
\[
x_A \ge 10
\]
\[
x_B \ge 15
\]

Minimum profit requirement:
\[
300x_A + 450x_B \ge 10000
\]

Nonnegativity:
\[
x_A \ge 0,\qquad x_B \ge 0.
\]

\section*{Notes on Implementation}

The code formulates and solves a linear program with objective sense ``maximize.'' The minimum profit target is imposed explicitly as a constraint in addition to maximizing profit. No integrality restrictions are used: \texttt{x\_A} and \texttt{x\_B} are continuous \texttt{LpVariable}s with lower bound $0$. The model is solved through PuLP using \texttt{GUROBI\_CMD(msg=0)}. No overtime decision, penalty term, or multiobjective/goal-programming structure is implemented; the capacities are enforced exactly as upper bounds read from the CSV data.
